of @@M0@@'';
to define the @@M1@@-th Chern class, we shall proceed analogously, bringing @@M2@@ to rank @@M3@@ by considering @@M4@@, and we shall take the alternating sum of the exterior powers of the element thus obtained; the following proposition is to be ``interpreted'' by setting @@M5@@:

\sgaStableTarget{sga6.e05.prop.3-2}\textbf{Proposition 3.2.} \uline{Let} @@M6@@ \uline{be a} @@M7@@-\uline{ring}.
\begin{enumerate}[label=\roman*)]
\item @@M8@@
\item \uline{Putting} @@M9@@, \uline{and} @@M10@@, \uline{we have}:
@@M11@@
@@M12@@
\end{enumerate}

We therefore see that knowing the operations @@M13@@ is equivalent to knowing the operations @@M14@@.

\uline{Proof}.

i) We have
@@M15@@
From @@M16@@, we obtain: @@M17@@, hence: @@M18@@, which gives the required relation.

ii) @@M19@@
We expand @@M20@@. @@M21@@; the coefficient of @@M22@@ is therefore:
@@M23@@

